La desinformación financiera en Colombia constituye una amenaza creciente que afecta la estabilidad del sistema financiero y la toma de decisiones de millones de ciudadanos. Según datos de Activa Research y WIN (2022), el 53\% de los colombianos enfrenta noticias falsas diariamente, mientras que el Foro Económico Mundial clasificó la desinformación como el riesgo global número uno a corto plazo en su informe de 2024. En el dominio financiero, esta problemática se manifiesta en rumores sobre quiebras bancarias, fraudes disfrazados de noticias institucionales y datos falsos sobre tasas de interés.

El presente proyecto propone el diseño, implementación y evaluación de un sistema híbrido que combina modelos de lenguaje de gran escala (LLMs) con un agente verificador de fuentes institucionales colombianas para la detección automatizada de desinformación financiera. El sistema opera en dos niveles: primero, un LLM clasifica la noticia en tres categorías (verdadera, dudosa o falsa) mediante prompting especializado; segundo, cuando el modelo presenta baja confianza en su clasificación (umbral inferior a 0.80), un agente inteligente basado en el framework ReAct consulta automáticamente fuentes como la Superintendencia Financiera de Colombia, el Banco de la República, el DANE y ColombiaCheck para emitir un veredicto fundamentado y trazable.

La investigación incluye la construcción del primer benchmark de evaluación de desinformación financiera colombiana (800--1.000 noticias etiquetadas), la comparación sistemática de múltiples LLMs en configuraciones zero-shot, few-shot y chain-of-thought, y la evaluación del impacto del agente validador sobre la precisión y explicabilidad del sistema. El proyecto se desarrolla en un periodo de ocho meses y se ejecuta sobre infraestructura accesible (Google Colab y APIs comerciales).

\textbf{Palabras clave}: desinformación financiera, modelos de lenguaje, agentes inteligentes, verificación de hechos, Colombia, procesamiento de lenguaje natural.
